\documentclass[11pt,a4paper]{article}
\usepackage[margin=25mm]{geometry}
\usepackage{amsmath,fontspec,unicode-math,array,longtable}
\setmainfont{Nimbus Roman}
\setmonofont{Nimbus Mono PS}
\setmathfont{MTPro2Math.otf}[Path=../out/]
% Use upright letters from the active math font for named operators.
\setoperatorfont{\symup}
% Named versions are used only to print alternatives on successive rows.
% A user's document normally selects its features once with setmathfont.
\setmathfont{MTPro2Math.otf}[Path=../out/,version=swash,
  RawFeature={+cv01}]
\setmathfont{MTPro2Math.otf}[Path=../out/,version=greekalt,
  RawFeature={+cv02,+cv03}]
\setmathfont{MTPro2Math.otf}[Path=../out/,version=betaonly,
  RawFeature={+cv02}]
\setmathfont{MTPro2Math.otf}[Path=../out/,version=deltaonly,
  RawFeature={+cv03}]
\setmathfont{MTPro2Math.otf}[Path=../out/,version=planckalt,
  RawFeature={+cv04}]
\setmathfont{MTPro2Math.otf}[Path=../out/,version=curly,
  RawFeature={+ss01}]
\setmathfont{MTPro2Math.otf}[Path=../out/,version=scriptalt,
  RawFeature={+ss02}]
\setmathfont{MTPro2Math.otf}[Path=../out/,version=curlytwo,
  RawFeature={+ss03}]
\setmathfont{MTPro2Math.otf}[Path=../out/,version=curlyalt,
  RawFeature={+ss01,+ss02}]
\setmathfont{MTPro2Math.otf}[Path=../out/,version=frakturalt,
  RawFeature={+ss04}]
\setmathfont{MTPro2Math.otf}[Path=../out/,version=heavy,
  RawFeature={+ss05}]
\setmathfont{MTPro2Math.otf}[Path=../out/,version=holey,
  RawFeature={+ss06}]
\setmathfont{MTPro2Math.otf}[Path=../out/,version=heavyholey,
  RawFeature={+ss05,+ss06}]
\setmathfont{MTPro2Math.otf}[Path=../out/,version=andalt,
  RawFeature={+cv05}]
\setmathfont{MTPro2Math.otf}[Path=../out/,version=slantedops,
  RawFeature={+ss07}]
\setmathfont{MTPro2Math.otf}[Path=../out/,version=straight,
  RawFeature={+cv06}]
\setmathfont{MTPro2Math.otf}[Path=../out/,version=largeparen,
  RawFeature={+cv07}]
\setmathfont{MTPro2Math.otf}[Path=../out/,version=largecurly,
  RawFeature={+cv08}]
\setmathfont{MTPro2Math.otf}[Path=../out/,version=plainhbrace,
  RawFeature={cv09=1}]
\setmathfont{MTPro2Math.otf}[Path=../out/,version=curlyhbrace,
  RawFeature={cv09=2}]
\setmathfont{MTPro2Math.otf}[Path=../out/,version=subsetalt,
  RawFeature={+cv10}]
\setmathfont{MTPro2Math.otf}[Path=../out/,version=orderalt,
  RawFeature={+cv11}]
\setmathfont{MTPro2Math.otf}[Path=../out/,version=negone,
  RawFeature={cv12=1,cv13=1}]
\setmathfont{MTPro2Math.otf}[Path=../out/,version=negtwo,
  RawFeature={cv12=2,cv13=2}]
\setmathfont{MTPro2Math.otf}[Path=../out/,version=negmixed,
  RawFeature={cv12=1,cv13=2}]
\setmathfont{MTPro2Math.otf}[Path=../out/,version=nsimalt,
  RawFeature={+cv14}]
\setmathfont{MTPro2Math.otf}[Path=../out/,version=curverel,
  RawFeature={+cv15,+cv16}]
\setmathfont{MTPro2Math.otf}[Path=../out/,version=proptoonly,
  RawFeature={+cv15}]
\setmathfont{MTPro2Math.otf}[Path=../out/,version=smileonly,
  RawFeature={+cv16}]
\setmathfont{MTPro2Math.otf}[Path=../out/,version=shortrel,
  RawFeature={+ss08}]
\setmathfont{MTPro2Math.otf}[Path=../out/,version=thickrel,
  RawFeature={+ss09}]
\setmathfont{MTPro2Math.otf}[Path=../out/,version=emptyalt,
  RawFeature={+cv17}]
\setmathfont{MTPro2Math.otf}[Path=../out/,version=marksalt,
  RawFeature={+cv18,+cv19}]
\setmathfont{MTPro2Math.otf}[Path=../out/,version=yenonly,
  RawFeature={+cv18}]
\setmathfont{MTPro2Math.otf}[Path=../out/,version=registeredonly,
  RawFeature={+cv19}]
\setmathfont{MTPro2Math.otf}[Path=../out/,version=cardsalt,
  RawFeature={+ss10}]
\usepackage[hidelinks,pdftitle={MTPro2 Math: Font guide},
  pdfauthor={}]{hyperref}
\setlength{\parindent}{0pt}
\setlength{\parskip}{5pt}
\setlength{\tabcolsep}{7pt}
\renewcommand{\arraystretch}{1.28}
\emergencystretch=2em
\clubpenalty=10000
\widowpenalty=10000
\newcommand{\codesize}{\fontsize{10.5pt}{12pt}\selectfont}
\newcommand{\code}[1]{{\codesize\texttt{\detokenize{#1}}}}
\AddToHook{env/verbatim/begin}{\codesize}
\newcommand{\sample}[2]{\begingroup\mathversion{#1}$#2$\endgroup}
\newcommand{\bbline}{\symbb{ABGMXYZ\;abcxyz\;0123456789}}
\newcommand{\bbgreek}{\BbbGamma\;\BbbPi\;\Bbbgamma\;\Bbbpi}
\newcommand{\bbitalic}{\mitBbbD\;\mitBbbd\;\mitBbbe\;\mitBbbi\;\mitBbbj}
\newcommand{\scriptline}{\symscr{CGILQSYZrz}}
\newcommand{\boldscriptline}{\symbfscr{CGILQSYZrz}}
\newcommand{\matrixsample}{\left(\begin{matrix}
  2&-1&0&0\\
  -1&2&-1&0\\
  0&-1&2&-1\\
  0&0&-1&2
\end{matrix}\right)}
\newcommand{\tentfunction}{\Lambda(x)=\left\{\begin{aligned}
  1+x,&\quad -1\leq x<0,\\
  1-x,&\quad 0\leq x\leq1,\\
  0,&\quad \text{otherwise}.
\end{aligned}\right.}
% Helpers below arrange specimens; they are not a font interface.
% U+00AE has no named unicode-math command; select the math family explicitly.
\newcommand{\registeredmark}{\Umathchar 0 \symoperators "00AE\relax}
\newcommand{\smallbrace}{\left\{\frac{1}{n}:n\in\symbb{N}\right\}}
\newcommand{\bracepair}{\overbrace{a+b+c}^{s}\quad\underbrace{x+y+z}_{t}}
\newcommand{\operatorsample}{\displaystyle
  \sum_{k=1}^{n}a_k\quad\prod_{k=1}^{n}b_k\quad\coprod_{i\in I}X_i}
\newcommand{\andsample}{p\wedge q\quad\displaystyle\bigwedge_{i=1}^{n}P_i}
\newcommand{\negatedsample}{a\nleq b\quad a\ngeq b\quad
  A\nsubseteq B\quad A\nsupseteq B}

\begin{document}

{\LARGE MTPro2 Math}\par
{\large Font guide}\hfill v1.5\par
\medskip
Use the generated font with LuaLaTeX and \code{unicode-math}:
\begin{verbatim}
\usepackage{unicode-math}
\setmathfont{MTPro2Math.otf}
\end{verbatim}
Use \code{MTPro2MathLite.otf} for Lite. The examples in this guide use Full.
No \code{mtpro2.sty} is needed. Building the font, installing the TTF on
Windows and using it in Word are described in \code{mtp2otf.pdf}.

\section{Letters and Greek}
Ordinary letters are mathematical italic. Use \code{\symup},
\code{\symit}, \code{\symbfup} and \code{\symbfit} to select a
mathematical alphabet explicitly:
\begin{center}
\begin{tabular}{@{}ll@{}}
\code{\symup{ABGMXYZabcxyz}} & $\symup{ABGMXYZabcxyz}$\\
\code{\symit{ABGMXYZabcxyz}} & $\symit{ABGMXYZabcxyz}$\\
\code{\symbfup{ABGMXYZabcxyz}} & $\symbfup{ABGMXYZabcxyz}$\\
\code{\symbfit{ABGMXYZabcxyz}} & $\symbfit{ABGMXYZabcxyz}$\\
\end{tabular}
\end{center}
The \code{\sym...} commands use the math font rather than the document's
text font. Bold italic requires Full. For Greek, the same selectors give
$\symup{\alpha\beta\Gamma\Omega}$,
$\symit{\alpha\beta\Gamma\Omega}$,
$\symbfup{\alpha\beta\Gamma\Omega}$ and
$\symbfit{\alpha\beta\Gamma\Omega}$.

Character variants (cv) select local letter or symbol forms; stylistic
sets (ss) select related designs across a group of characters. Both can
be selected for the whole document.
For example, swash z is selected in the preamble:
\begin{verbatim}
\setmathfont{MTPro2Math.otf}[RawFeature={+cv01}]
\end{verbatim}
\begin{center}
\begin{tabular}{@{}lll@{}}
Default & \sample{normal}{z\quad\symbfit{z}} & swash z\\
\code{cv01} & \sample{swash}{z\quad\symbfit{z}} & \\[5pt]
Default & \sample{normal}{\symup{\beta\delta}\quad\symit{\beta\delta}\quad
  \symbfup{\beta\delta}\quad\symbfit{\beta\delta}} & beta, delta\\
\code{cv02,cv03} & \sample{greekalt}{\symup{\beta\delta}\quad\symit{\beta\delta}\quad
  \symbfup{\beta\delta}\quad\symbfit{\beta\delta}} & \\[5pt]
Default & \sample{normal}{\hbar} & Planck sign\\
\code{cv04} & \sample{planckalt}{\hbar} & \\
\end{tabular}
\end{center}
These choices are available in both editions; their bold counterparts
require Full. Under \code{unicode-math}, \code{\hbar} and
\code{\hslash} name the same character, whose design is selected by
\code{cv04}.

Beta and delta are independent: \code{cv02} changes beta, and
\code{cv03} changes delta. The row above enables both.
Features can be combined: \code{RawFeature={+cv01,+cv02}} selects swash z
and alternate beta while leaving delta unchanged.
For a feature with numbered choices, use, for example,
\code{RawFeature={cv09=2}}. Omit a feature to keep its default design.

\newpage
\section{Script and Fraktur}
These alphabets are available in Full. Use \code{\symscr} and \code{\symbfscr} for regular and bold Script.
\code{ss01} selects Curly Script. \code{ss02} changes selected letters
in regular and bold Script, including alternatives in the Curly family.
\code{ss03} selects the alternate Curly series directly.

\begin{center}
\code{\symscr{CGILQSYZrz}}\par\smallskip
\begin{tabular}{@{}ll@{}}
Default & \sample{normal}{\scriptline}\\
\code{ss01} & \sample{curly}{\scriptline}\\
\code{ss02} & \sample{scriptalt}{\scriptline}\\
\code{ss03} & \sample{curlytwo}{\scriptline}\\
\code{ss01,ss02} & \sample{curlyalt}{\scriptline}\\
\end{tabular}
\end{center}
\begin{verbatim}
\setmathfont{MTPro2Math.otf}[RawFeature={+ss01,+ss02}]
\end{verbatim}
\code{ss03} takes precedence over the other regular Script choices.
It matches \code{ss01,ss02} for regular Script, but does not itself enable
bold Script alternatives. Curly G/M/N/Q/Y also have alternatives under
\code{ss02}.
The Curly selectors affect regular Script; bold Script retains its own
design, with its alternatives controlled by \code{ss02}.

\begin{center}
\code{\symbfscr{CGILQSYZrz}}\par\smallskip
\begin{tabular}{@{}ll@{}}
Default & \sample{normal}{\boldscriptline}\\
\code{ss02} & \sample{scriptalt}{\boldscriptline}\\
\end{tabular}
\end{center}

Fraktur is selected with \code{\symfrak} or \code{\symbffrak}.
Its alternatives use \code{ss04}, independently of the Script choices.

\begin{center}
\begin{tabular}{@{}lll@{}}
 & \code{\symfrak{IRYxy}} & \code{\symbffrak{IRYxy}}\\
Default & \sample{normal}{\symfrak{IRYxy}}
  & \sample{normal}{\symbffrak{IRYxy}}\\
\code{ss04} & \sample{frakturalt}{\symfrak{IRYxy}}
  & \sample{frakturalt}{\symbffrak{IRYxy}}\\
\end{tabular}
\end{center}
\begin{verbatim}
\setmathfont{MTPro2Math.otf}[RawFeature={+ss02,+ss04}]
\end{verbatim}

\newpage
\section{Blackboard}
Blackboard is available in Full. Use \code{\symbb{...}} for double-struck letters and digits.
The default design and three upright alternatives are shown below.

\begin{center}
\code{\symbb{ABGMXYZ\;abcxyz\;0123456789}}\par\smallskip
\begin{tabular}{@{}ll@{}}
Default & \sample{normal}{\bbline}\\
\code{ss05} & \sample{heavy}{\bbline}\\
\code{ss06} & \sample{holey}{\bbline}\\
\code{ss05,ss06} & \sample{heavyholey}{\bbline}\\
\end{tabular}
\end{center}
\code{ss05} selects heavier upright letters. \code{ss06} selects the
Holey Roman family; combine them for its heavier upright design.
For example:
\begin{verbatim}
\setmathfont{MTPro2Math.otf}[RawFeature={+ss05,+ss06}]
\end{verbatim}

Four Greek characters and five italic characters have standard Unicode
inputs. \code{ss06} changes these to the corresponding Holey designs;
\code{ss05} leaves them unchanged.

\begin{center}
\begin{tabular}{@{}lll@{}}
 & Greek & Italic\\
Default & \sample{normal}{\bbgreek} & \sample{normal}{\bbitalic}\\
\code{ss06} & \sample{holey}{\bbgreek} & \sample{holey}{\bbitalic}\\
\end{tabular}
\end{center}
\begin{verbatim}
\BbbGamma \BbbPi \Bbbgamma \Bbbpi
\mitBbbD \mitBbbd \mitBbbe \mitBbbi \mitBbbj
\end{verbatim}
The italic Blackboard selection is limited to these five characters.

\newpage
\section{Operators and formula sizes}
To use the math font's upright letters for named operators such as
\code{\sin} and \code{\cos}, set:
\begin{verbatim}
\setoperatorfont{\symup}
\end{verbatim}
Large operators have ordinary text and display designs. The math
style selects them automatically; \code{\limits} and \code{\nolimits} control the
placement of limits.
\[
  \textstyle\sum_{k=1}^{n}a_k\quad\int_0^1 f(x)\,dx
  \qquad\displaystyle\sum_{k=1}^{n}a_k\quad\int_0^1 f(x)\,dx
\]
\code{cv05} changes the three optical designs of \code{\wedge}
and the large operator \code{\bigwedge} at every supported size:
\begin{center}
\begin{tabular}{@{}ll@{}}
Default & \sample{normal}{\andsample}\\[9pt]
\code{cv05} & \sample{andalt}{\andsample}\\
\end{tabular}
\end{center}
\code{ss07} selects slanted sums, products and coproducts at every supported size.
\begin{center}
\begin{tabular}{@{}ll@{}}
Default & \sample{normal}{\operatorsample}\\[9pt]
\code{ss07} & \sample{slantedops}{\operatorsample}\\
\end{tabular}
\end{center}
\begin{verbatim}
\setmathfont{MTPro2Math.otf}[RawFeature={+ss07}]
\end{verbatim}
Both operator features are available in Full and Lite. The features select
design families; the engine's size request selects a stage within a family.
Other operators retain their own designs:
\[
  \bigcap_{i=1}^{n}A_i\qquad\bigcup_{i=1}^{n}A_i
  \qquad\bigotimes_{i=1}^{n}V_i\qquad\oint_C\vec{F}\cdot d\vec{r}
\]
The original package's extra-large operator commands are not provided.
Use a size request as described below instead.

\newpage
\subsection*{Requesting a larger operator}
With ordinary \code{unicode-math} input, a tall fraction after \code{\int}
does not enlarge the integral sign. To select a larger design in LuaLaTeX,
set \code{\Umathoperatorsize} for the required math style:
\begin{verbatim}
\begingroup
\Umathoperatorsize\displaystyle=3em
\[
  \int_0^1 \frac{1}{1+\frac{x}{1+x^2}}\,dx
\]
\endgroup
\end{verbatim}
\begin{center}
\begin{tabular}{@{}cc@{}}
Ordinary display & Requested size \code{3em} \\[10pt]
$\displaystyle\int_0^1\frac{1}{1+\frac{x}{1+x^2}}\,dx$
& \begingroup\Umathoperatorsize\displaystyle=3em
$\displaystyle\int_0^1\frac{1}{1+\frac{x}{1+x^2}}\,dx$\endgroup
\\
\end{tabular}
\end{center}
The value is a size request, not a scale factor: the engine chooses from
the font's ready-made designs. It does not continuously stretch the glyph.
The group limits the setting to this example. Within the group, the setting
also applies to other operators in display style, such as \code{\sum} and
\code{\prod}. Stylistic sets such as \code{ss07} choose the design
family independently of the size request.

Larger designs are available in both Full and Lite. The supported series
and their number of sizes are listed in \code{mtp2otf.pdf}, Appendix B.
The command above is specific to LuaTeX; Word controls operator growth
through its own equation settings.

\subsection*{Scripts, fractions and radicals}
Subscripts and superscripts select the source's optical sizes automatically:
\[
  x+\alpha\qquad 2^{x+\alpha}\qquad 2^{2^{x+\alpha}}
\]
Radicals and fractions use the ordinary commands:
\begin{verbatim}
\sqrt{\frac{a^2+b^2}{c^2+d^2}} \qquad \sqrt[3]{x+y}
\end{verbatim}
\[
  \sqrt{\frac{a^2+b^2}{c^2+d^2}}\qquad\sqrt[3]{x+y}
\]

\newpage
\section{Parentheses and braces}
Use \code{\left} and \code{\right} to size delimiters to a formula.
The default series uses the original ordinary growth range, followed by
an assembly of repeatable parts. The default braces use the morphed
design, which approaches the assembled shape as it grows.

\code{cv07} extends parentheses through the large ready-made designs.
The two examples below use the same input and font size.

\begin{center}
\begin{tabular}{@{}ll@{}}
Default & \sample{normal}{\matrixsample}\\[10pt]
\code{cv07} & \sample{largeparen}{\matrixsample}\\
\end{tabular}
\end{center}
\begin{verbatim}
\setmathfont{MTPro2Math.otf}[RawFeature={+cv07}]
\end{verbatim}

For braces, \code{cv06} selects straight braces with the ordinary growth
range. \code{cv08} selects Curly braces with an extended ready-made range.

At larger sizes the default and straight series use the same assembly.
Each row shows a small set and a larger piecewise function:
\begin{center}
\begin{tabular}{@{}lll@{}}
Default & \sample{normal}{\smallbrace}
  & \sample{normal}{\tentfunction}\\[10pt]
\code{cv06} & \sample{straight}{\smallbrace}
  & \sample{straight}{\tentfunction}\\[10pt]
\code{cv08} & \sample{largecurly}{\smallbrace}
  & \sample{largecurly}{\tentfunction}\\
\end{tabular}
\end{center}
\begin{verbatim}
\setmathfont{MTPro2Math.otf}[RawFeature={+cv07,+cv08}]
\end{verbatim}
Use paired delimiters, or end with \code{\right.} for a one-sided brace:
\begin{verbatim}
\left( ... \right)     \left\{ ... \right\}
\left\{ ... \right.
\end{verbatim}

The large series reaches 28.8\,em before continuing with the ordinary
assembly. \code{cv06} takes precedence over \code{cv08}; choose one brace
style. Parenthesis selection remains independent.

\newpage
\section{Accents and horizontal braces}
Use a fixed accent for a letter and a wide accent for a group. These
constructions are available in both editions.
\begin{center}
\begin{tabular}{@{}ll@{}}
\code{\hat{x}, \widehat{ABC}} & $\hat{x}\quad\widehat{ABC}$\\[2pt]
\code{\tilde{x}, \widetilde{ABC}} & $\tilde{x}\quad\widetilde{ABC}$\\[2pt]
\code{\check{x}, \widecheck{ABC}} & $\check{x}\quad\widecheck{ABC}$\\[2pt]
\code{\bar{M}, \wideoverbar{M+N}} & $\bar{M}\quad\wideoverbar{M+N}$\\[2pt]
\code{\vec{v}, \overrightarrow{AB}} & $\vec{v}\quad\overrightarrow{AB}$\\[2pt]
\code{\overleftarrow{AB}} & $\overleftarrow{AB}$\\[2pt]
\code{\dot{x}, \ddot{x}, \hat{i}} & $\dot{x}\quad\ddot{x}\quad\hat{i}$\\
\end{tabular}
\end{center}
\code{\bar} and \code{\overbar} select the fixed accent;
\code{\wideoverbar} selects the stretching accent.
\code{\overline} draws a rule instead. An operator such as
$\hat{H}$ uses the ordinary \code{\hat{H}}.

Horizontal parentheses and braces cover their contents. For example:
\begin{verbatim}
\overparen{AB} \qquad \underbrace{V\times\cdots\times V}_{n}
\end{verbatim}
\[
  \overparen{AB}\qquad\underbrace{V\times\cdots\times V}_{n}
\]
\code{cv09} selects the design of both overbraces and underbraces:
\begin{center}
\begin{tabular}{@{}ll@{}}
Default & \sample{normal}{\bracepair}\\[10pt]
\code{cv09=1}: plain & \sample{plainhbrace}{\bracepair}\\[10pt]
\code{cv09=2}: Curly & \sample{curlyhbrace}{\bracepair}\\
\end{tabular}
\end{center}
\begin{verbatim}
\setmathfont{MTPro2Math.otf}[RawFeature={cv09=2}]
\end{verbatim}
This choice is independent of the vertical-brace features. Use:
\begin{verbatim}
\overbrace{...}^{...}     \underbrace{...}_{...}
\end{verbatim}

\newpage
\section{Relations}
The following Full-edition features change relation designs without
changing their input commands or mathematical meanings.
\begin{center}
\begin{tabular}{@{}lll@{}}
Default & \sample{normal}{\subsetneq\;\supsetneq\;\subsetneqq\;\supsetneqq} & \\
\code{cv10} & \sample{subsetalt}{\subsetneq\;\supsetneq\;\subsetneqq\;\supsetneqq}
  & \code{\subsetneq \supsetneq}\\
 & & \code{\subsetneqq \supsetneqq}\\[5pt]
Default & \sample{normal}{\lneqq\;\gneqq} & \\
\code{cv11} & \sample{orderalt}{\lneqq\;\gneqq}
  & \code{\lneqq \gneqq}\\[5pt]
\end{tabular}
\end{center}
Negated comparisons use \code{cv12}; negated inclusions use \code{cv13}.
Each independently accepts 0 (off), 1 (first alternative) or 2 (double
equality). Value 1 selects slanted equality for comparisons and the AMS
alternate for inclusions. Use their direct commands:
\begin{verbatim}
a\nleq b \quad a\ngeq b \quad A\nsubseteq B \quad A\nsupseteq B
\end{verbatim}
\begin{center}
\begin{tabular}{@{}ll@{}}
Default & \sample{normal}{\negatedsample}\\[4pt]
\code{cv12=1,cv13=1} & \sample{negone}{\negatedsample}\\[4pt]
\code{cv12=2,cv13=2} & \sample{negtwo}{\negatedsample}\\
\end{tabular}
\end{center}
The not-similar sign is independent:
\begin{center}
\begin{tabular}{@{}lll@{}}
Default & \sample{normal}{a\nsim b} & \code{a\nsim b}\\
\code{cv14} & \sample{nsimalt}{a\nsim b} & \\
\end{tabular}
\end{center}
\begin{verbatim}
\setmathfont{MTPro2Math.otf}[RawFeature={cv12=1,cv13=1,+cv14}]
\end{verbatim}

\code{cv15} selects the proportional-to form; \code{cv16} selects smile
and frown. They can be enabled independently.
\begin{center}
\begin{tabular}{@{}lll@{}}
Default & \sample{normal}{\propto\;\smile\;\frown} & \\
\code{cv15,cv16} & \sample{curverel}{\propto\;\smile\;\frown}
  & \code{\propto \smile \frown}\\
\end{tabular}
\end{center}
\code{ss08} selects shorter relation bars; it does not change the
delimiter bars in \code{\left|x\right|} or
\code{\left\Vert v\right\Vert}. \code{ss09} selects thicker
similarity and approximation signs.
\begin{center}
\begin{tabular}{@{}lll@{}}
Default & \sample{normal}{\mid\;\nmid\;\parallel\;\nparallel} & \\
\code{ss08} & \sample{shortrel}{\mid\;\nmid\;\parallel\;\nparallel}
  & \code{\mid \nmid \parallel \nparallel}\\[5pt]
Default & \sample{normal}{\sim\;\approx} & \\
\code{ss09} & \sample{thickrel}{\sim\;\approx}
  & \code{\sim \approx}\\
\end{tabular}
\end{center}

\newpage
\section{Other symbols}
\code{cv17} selects the alternative empty-set sign in Full.
\code{\emptyset} and \code{\varnothing} refer to the same Unicode
character under \code{unicode-math}.
\begin{center}
\begin{tabular}{@{}lll@{}}
Default & \sample{normal}{\emptyset} & \code{\emptyset}\\
\code{cv17} & \sample{emptyalt}{\emptyset} & \\[5pt]
Default & \sample{normal}{\mathyen\;\registeredmark}
  & \code{\mathyen}, U+00AE\\
\code{cv18,cv19} & \sample{marksalt}{\mathyen\;\registeredmark} & \\[5pt]
Default & \sample{normal}{\spadesuit\;\clubsuit}
  & \code{\spadesuit \clubsuit}\\
\code{ss10} & \sample{cardsalt}{\spadesuit\;\clubsuit} & \\
\end{tabular}
\end{center}
\code{cv18} changes yen and \code{cv19} changes the registered sign in Full;
they are independent. \code{ss10} gives shaded spades and clubs in both editions.

The usual symbol commands include \code{\ell}, \code{\partial},
\code{\nabla}, \code{\infty}, \code{\dagger} and \code{\ddagger}:
\[
  \ell\qquad\partial\qquad\nabla\qquad\infty\qquad\dagger\qquad\ddagger
\]
Mathematical sans-serif and monospace alphabets, a complete bold-symbol
version, and a compatibility package for MTPro2-specific macros are
not included.

\newpage
\section{Combining features}
These features can be selected independently.
Each example keeps the other symbol at its default unless explicitly enabled.
\begin{center}
\begin{tabular}{@{}ll@{}}
Default & \sample{normal}{\beta+\delta}\\
\code{cv02} & \sample{betaonly}{\beta+\delta}\\
\code{cv03} & \sample{deltaonly}{\beta+\delta}\\[6pt]
Default & \sample{normal}{\negatedsample}\\
\code{cv12=1,cv13=2} & \sample{negmixed}{\negatedsample}\\[6pt]
Default & \sample{normal}{a\propto b\quad x\smile y\quad x\frown y}\\
\code{cv15} & \sample{proptoonly}{a\propto b\quad x\smile y\quad x\frown y}\\
\code{cv16} & \sample{smileonly}{a\propto b\quad x\smile y\quad x\frown y}\\[6pt]
Default & \sample{normal}{\mathyen\quad\registeredmark}\\
\code{cv18} & \sample{yenonly}{\mathyen\quad\registeredmark}\\
\code{cv19} & \sample{registeredonly}{\mathyen\quad\registeredmark}\\[6pt]
\end{tabular}
\end{center}
\begin{verbatim}
\setmathfont{MTPro2Math.otf}[RawFeature={cv12=1,cv13=2}]
\end{verbatim}

\newpage
\section{Three math sizes}
Each group repeats an expression at normal size, in a superscript and in
an exponent of a superscript. The exponent construction is used here
only to compare the three optical sizes.

\[
  A,\ldots,Z\quad a,\ldots,z\quad\Gamma,\Omega\quad\alpha_i,\omega^j
\]
\[
  2^{A,\ldots,Z\quad a,\ldots,z\quad\Gamma,\Omega\quad\alpha_i,\omega^j}
\]
\[
  2^{2^{A,\ldots,Z\quad a,\ldots,z\quad\Gamma,\Omega\quad\alpha_i,\omega^j}}
\]

\[
  (f\circ g)'(x)=f'(g(x))\,g'(x)
\]
\[
  2^{(f\circ g)'(x)=f'(g(x))\,g'(x)}
\]
\[
  2^{2^{(f\circ g)'(x)=f'(g(x))\,g'(x)}}
\]

\[
  \sum_{k=1}^{n}\frac{1}{k(k+1)}=\frac{n}{n+1}
\]
\[
  2^{\sum_{k=1}^{n}\frac{1}{k(k+1)}=\frac{n}{n+1}}
\]
\[
  2^{2^{\sum_{k=1}^{n}\frac{1}{k(k+1)}=\frac{n}{n+1}}}
\]

\[
  \int_{-\infty}^{\infty}e^{-x^2}\,dx=\sqrt{\pi}
\]
\[
  2^{\int_{-\infty}^{\infty}e^{-x^2}\,dx=\sqrt{\pi}}
\]
\[
  2^{2^{\int_{-\infty}^{\infty}e^{-x^2}\,dx=\sqrt{\pi}}}
\]

\[
  \omega=v\,dx+w\,dy,\qquad d\omega=\left(\frac{\partial w}{\partial x}-\frac{\partial v}{\partial y}\right)\,dx\wedge dy
\]
\[
  2^{\omega=v\,dx+w\,dy,\qquad d\omega=\left(\frac{\partial w}{\partial x}-\frac{\partial v}{\partial y}\right)\,dx\wedge dy}
\]
\[
  2^{2^{\omega=v\,dx+w\,dy,\qquad d\omega=\left(\frac{\partial w}{\partial x}-\frac{\partial v}{\partial y}\right)\,dx\wedge dy}}
\]

\[
  R_{ijkl}=-R_{jikl}=-R_{ijlk}=R_{klij}
\]
\[
  2^{R_{ijkl}=-R_{jikl}=-R_{ijlk}=R_{klij}}
\]
\[
  2^{2^{R_{ijkl}=-R_{jikl}=-R_{ijlk}=R_{klij}}}
\]

\end{document}
